\section{Discussion}

\paragraph{Adaptation should pay for the decisions it changes.}
FLEX-UF identifies a budget-dependent role for spatial allocation.
Under source calibration, tight quality targets leave room for useful
placement. Near the shallowest-exit ceiling, a blind mixture reaches
nearly the same arithmetic saving. This motivates a budget-conditioned choice between fixed, blind and
learned policies, with its switching rule fixed on validation data.

\paragraph{Reuse and specialisation are different commitments.}
Shared exits retain one representation and the work already performed;
independent experts can adapt their encoders and synthesis capacity more
freely. The bank adds model residency, entropy boundaries and occupancy-dependent
execution costs. Its quality must be paired with those costs, not with
shared-prefix arithmetic. Multiple exits inside each expert are outside
the present design.

\paragraph{The remaining evidence gap.}
The source-calibrated study measures available allocation value, while
the fixed-control replay leaves the routing increment unresolved.
Content-grouped calibration on untouched images is the next test of that
increment. Native bitstreams and paired wall-clock measurements must then
include entropy work, prediction, tile movement, grouping and fallback.
The current confidence intervals cover sample variation at fixed weights,
not independent training seeds. Results are confined to intra reconstruction:
YCbCr-MSE and RGB diagnostics do not establish perceptual equivalence,
temporal consistency or performance of DCVC-UF's chunk-based video path.

\section{Conclusion}
FLEX-UF makes synthesis depth a spatial decision inside DCVC-UF.
A common latent and prefix support aligned early exits, while a shrinking
tile batch turns stopping decisions into omitted block evaluations.
The measured arithmetic savings have two sources: cheaper average depth
and a smaller, budget-dependent increment from learned placement.
Blind mixtures and fixed-control replay show why these must be evaluated
separately. The independent-depth study extends the allocation question
to codec selection; matched training and complete execution measurements
will determine when specialisation justifies giving up shared computation.
